\section{State of the art}
\label{sec:related}

\begin{inbrief}
The problem sits where six lines of work meet: zonal architectures,
task allocation, task-chain analysis and LET, TSN analysis, shared-resource
contention, and system-level analysis and the joint design of computation
and communication.
\Cref{tab:related} places the closest works against the features of the
problem. None covers allocation, priorities, task-chains, LET and network
bounds together on a zonal platform, and none sizes the cores.
\end{inbrief}

\subsection{Zonal and centralized E/E architectures}

Bandur et al.\ \cite{bandur2021centralized} make the case for centralized
automotive architectures: fewer, more powerful computers, a service-based
software platform and a backbone network. Askaripoor et al.\
\cite{askaripoor2022synthesis} survey E/E architecture synthesis, of which
the allocation of software to computing units is one step. Mauser and
Wagner \cite{mauser2025centralization} combine a literature review and
industrial interviews to identify seven system properties that limit
current E/E architectures (bus load, functional safety and computing
power among them) and use them to assess the potential of
centralization.

Quantitative studies of zonal platforms have so far looked at the network
and the harness. Park et al.\ \cite{park2024zonal} simulate the end-to-end delay of TSN
streams for 2 to 10 zones and estimate the length and weight of the
harness. Aghababaiyan et
al.\ \cite{aghababaiyan2025deterministic} schedule tasks between zone ECUs
and a central unit over several topologies, and report deadline
satisfaction from simulation, without a worst-case analysis. Zaheri and
Reuss \cite{zaheri2025asil} allocate tasks to ECUs with a mixed-integer
program under ISO~26262 constraints \cite{iso26262} (ASIL decomposition,
memory) on a case of six tasks. None of these works sizes the cores of the
ECUs under a worst-case analysis of the task-chains and the network.

\subsection{Task allocation and partitioned scheduling}

Allocating periodic tasks to processors is a bin-packing problem, and the
decision version of bin packing is NP-complete in the strong sense
\cite{garey1979computers}. Partitioned scheduling therefore
relies on bin-packing heuristics (first, best or worst fit by decreasing
utilization) with a per-core schedulability test \cite{davis2011survey},
and task splitting extends them \cite{lakshmanan2009partitioned}. First fit decreasing uses at most $11/9$ of
the optimal number of bins plus $6/9$ \cite{dosa2007tight}. On
heterogeneous processors the size of an item depends on the bin
\cite{baruah2004heterogeneous}.

For distributed systems, Tindell et al.\ \cite{tindell1992allocating}
searched allocations by simulated annealing, with response-time analysis
of tasks and messages. Zheng et al.\ \cite{zheng2007definition} and Zhu et
al.\ \cite{zhu2012optimization} optimize the allocation and the priorities
together with mixed-integer programming, under end-to-end deadlines on a
CAN-based platform. Zhu et al.\ \cite{zhu2009optimizing} optimize the
\emph{extensibility} of such a system, the room left for WCET increases.
Cambazard et al.\ \cite{cambazard2004decomposition} solve a hard real-time
allocation problem exactly by decomposition: a master problem allocates,
and a schedulability subproblem returns explanations that cut the master;
\cref{sec:exact} follows the same idea. Davare et al.\
\cite{davare2007period} optimize periods for latency. Bhat et al.\
\cite{bhat2017practical} allocate primary and backup tasks for software
fault tolerance in automotive systems, so that replicas never share a
processor.

Within a processor, the WCRT of fixed-priority tasks is the least fixed
point of a recurrence \cite{joseph1986finding,audsley1993applying}.
Audsley's algorithm \cite{audsley2001priority} finds a feasible priority
order whenever one exists, for any test that does not depend on the
relative order of the higher-priority tasks \cite{davis2016review}. Davis and Burns
\cite{davis2007robust} refine it to maximize the tolerance to
interference, and Davis et al.\ \cite{davis2016review} review its
extensions. The margin of a schedulable system is measured by the
critical scaling factor \cite{lehoczky1989rate} and, more generally, by
sensitivity analysis \cite{bini2008sensitivity}. The priority orders of
tasks with precedence constraints are a restriction these methods do not
handle directly; \cref{sec:audsley} adapts Audsley's argument to it.

\subsection{Task-chains and logical execution time}

Feiertag et al.\ \cite{feiertag2008compositional} defined the reaction
time and the data age of task-chains whose tasks communicate through
registers. Becker et al.\ \cite{becker2017end} bound the maximum data age
of automotive task-chains for several levels of system knowledge,
including implicit and LET communication. Dürr et al.\
\cite{durr2019end} and Günzel et al.\ \cite{gunzel2021timing,
gunzel2023compositional} extend the bounds to sporadic and asynchronous
distributed task-chains, and Günzel et al.\ \cite{gunzel2023equivalence}
show that maximum reaction time and maximum data age are equal.

LET originates in Giotto \cite{henzinger2003giotto}; Kirsch and Sokolova
\cite{kirsch2012let} describe the paradigm. Bosch reported its industrial
use for multicore engine control \cite{hamann2017communication} and
published the statistics of an engine-management software as a WATERS
benchmark \cite{kramer2015waters}. Biondi and Di Natale
\cite{biondi2018achieving} implement LET on multicores and measure its
overhead. Martinez et al.\ \cite{martinez2018analytical,martinez2020end}
characterize end-to-end delays under LET and other communication models.
Pazzaglia et al.\ \cite{pazzaglia2019optimizing} optimize the deployment
of functions onto the tasks and cores of one multicore under LET, and
Paladino et al.\ \cite{paladino2024optimizing} assign per-core priorities
to minimize the latency of task-chains that span cores. Gemlau et al.\
\cite{gemlau2021system} extend LET to distributed systems by giving the
network its own logical intervals. All these works analyse or optimize
task-chains on a \emph{given} platform; the sizing of the platform is
outside their scope.

\subsection{Timing analysis of TSN}

Network Calculus \cite{cruz1991calculus,leboudec2001network} bounds delays
with arrival and service curves. For the credit-based shaper, Queck
\cite{queck2012analysis} and De Azua and Boyer \cite{azua2014complete}
derive service curves for the SR classes A and B; Mohammadpour et al.\
\cite{mohammadpour2018latency} and Zhao et al.\ \cite{zhao2021latency}
bound the credit of any number of classes and give the service curves
from which \cref{sec:nc} takes its rate-latency form. For time-aware shaping, Zhao et al.\
\cite{zhao2018worst} give Network Calculus bounds and Craciunas et al.\
\cite{craciunas2016scheduling} synthesize gate schedules. Thiele and Ernst
\cite{thiele2016formal} analyse frame preemption, and Zhao et al.\
\cite{zhao2022quantitative} compare the shapers.

Asynchronous traffic shaping comes from the urgency-based scheduler of
Specht and Samii \cite{specht2016urgency}. Le Boudec
\cite{leboudec2018theory} proved that interleaved regulators placed
behind a FIFO system do not increase its worst-case delay (``shaping for
free''), so that the bursts of a flow can be bounded at their source value
at every hop. Thomas et al.\ \cite{thomas2022worst} analyse packet
replication and elimination: the elimination increases the burstiness of
a flow, and an interleaved regulator placed after it loses the
shaping-for-free property. Most TSN analyses take the streams as given; in
this paper the streams depend on the mapping.

\subsection{Shared-resource contention}

On a multicore, tasks of different cores delay each other on shared
caches, buses and memory. Paolieri et al.\ \cite{paolieri2009hardware}
bound the delay of a request under round-robin arbitration by one request
of every other core, which makes WCETs composable. Davis et al.\
\cite{davis2018extensible} integrate such interference into
response-time analysis (multicore response time analysis). Allocation
studies usually assume WCETs that already include the worst
interference; \cref{sec:contention} makes the interference depend on the
mapping instead. Zhu et al.\ \cite{zhu2026reducing} refine the shared-cache
interference of cause-effect chains on a given mapping of chains to
cores, which lowers their end-to-end latency bounds.

\subsection{System-level analysis and joint design}

Compositional system-level analysis, as in SymTA/S
\cite{henia2005system}, propagates event models between the local
analyses of processors and buses, and bounds end-to-end delays for a
given mapping. Craciunas and Oliver \cite{craciunas2016combined}
synthesize task and network schedules together for time-triggered
systems. Pop et al.\ \cite{pop2016design} optimize the routing and
scheduling of time-triggered and AVB traffic in TSN. Both assume the tasks
are already allocated. Arestova et al.\ \cite{arestova2022itans} schedule
multi-rate cause-effect chains of tasks and of their TSN flows together,
by an incremental heuristic under end-to-end deadlines; they compute a
schedule, where this paper allocates the tasks and bounds the delays of
fixed-priority tasks and of shaped streams. Zhu et al.\ \cite{zhu2012optimization} do allocate tasks, to ECUs connected
by a CAN bus with fixed-priority messages, with a mixed-integer program.

\begin{table*}[t]
\centering\footnotesize
\caption{Positioning against the closest works. \ding{51}: covered;
$\circ$: partly, or under strong assumptions; --: not addressed (as
described in the publications). \emph{Alloc.}: allocation of tasks to cores
or ECUs; \emph{Prio.}: fixed-priority assignment; \emph{Chains}: end-to-end
bounds of multi-rate task-chains; \emph{LET}: LET communication;
\emph{Net.}: worst-case network bounds; \emph{Zonal}: zonal bindings and
spatial constraints (memory, replicas, co-location); \emph{Het.}:
heterogeneous core speeds or types; \emph{Size}: minimizes the number of
cores; \emph{Exact}: exact method or optimality gap measured;
\emph{Life}: evolution of the software after the design; \emph{Scale}:
largest reported case.}
\label{tab:related}
\setlength{\tabcolsep}{3.4pt}
\renewcommand{\arraystretch}{1.08}
\rowcolors{2}{brandlight}{white}
\begin{tabular}{@{}l*{10}{c}l@{}}
\toprule
\rowcolor{white}
work & Alloc. & Prio. & Chains & LET & Net. & Zonal & Het. & Size & Exact & Life & Scale \\
\midrule
Tindell et al.\ 1992 \cite{tindell1992allocating} & \ding{51} & \ding{51} & $\circ$ & -- & $\circ$ & -- & -- & -- & -- & -- & -- \\
Cambazard et al.\ 2004 \cite{cambazard2004decomposition} & \ding{51} & \ding{51} & -- & -- & $\circ$ & $\circ$ & -- & -- & \ding{51} & -- & -- \\
Zhu et al.\ 2009/12 \cite{zhu2009optimizing,zhu2012optimization} & \ding{51} & \ding{51} & \ding{51} & -- & \ding{51} (CAN) & -- & $\circ$ & -- & \ding{51} & $\circ$ & -- \\
Bhat et al.\ 2017 \cite{bhat2017practical} & \ding{51} & -- & -- & -- & -- & $\circ$ (replicas) & -- & $\circ$ & -- & -- & -- \\
Pazzaglia et al.\ 2019 \cite{pazzaglia2019optimizing} & \ding{51} & $\circ$ & \ding{51} & \ding{51} & -- & -- & -- & -- & $\circ$ & -- & one multicore \\
Paladino et al.\ 2024 \cite{paladino2024optimizing} & -- & \ding{51} & \ding{51} & $\circ$ & -- & -- & -- & -- & -- & -- & one multicore \\
Gemlau et al.\ 2021 \cite{gemlau2021system} & -- & -- & \ding{51} & \ding{51} & $\circ$ & -- & -- & -- & -- & -- & -- \\
Günzel et al.\ 2021/23 \cite{gunzel2021timing,gunzel2023compositional} & -- & -- & \ding{51} & \ding{51} & $\circ$ & -- & -- & -- & -- & -- & -- \\
Craciunas, Oliver 2016 \cite{craciunas2016combined} & -- & $\circ$ (TT) & \ding{51} & -- & \ding{51} (TT) & -- & -- & -- & \ding{51} & -- & -- \\
Zaheri, Reuss 2025 \cite{zaheri2025asil} & \ding{51} & $\circ$ & $\circ$ & -- & -- & $\circ$ (ASIL, memory) & \ding{51} & -- & \ding{51} & -- & 6 tasks \\
Aghababaiyan et al.\ 2025 \cite{aghababaiyan2025deterministic} & \ding{51} & -- & -- & -- & -- (simulated) & \ding{51} & \ding{51} & -- & -- & -- & -- \\
\midrule
\rowcolor{white}
\textbf{This paper} & \ding{51} & \ding{51} & \ding{51} & \ding{51} & \ding{51} (4 TSN) & \ding{51} & \ding{51} & \ding{51} & \ding{51} & \ding{51} & \textbf{13\,869 tasks} \\
\bottomrule
\end{tabular}
\end{table*}

\subsection{Positioning}

\Cref{tab:related} summarizes the gap. Works that allocate tasks either
ignore the network or assume a bus with fixed-priority messages; works
that bound task-chains, LET or TSN take the allocation as given;
contention is analysed on a given mapping; and no work both sizes the
cores of a zonal platform and bounds its task-chains with all these
couplings in one system-level analysis. The aim of this paper is a sound, explainable
and scalable method that a system architect can run on the software of a
whole vehicle in early design. It compares platforms and design rules
(core types, network rate, TSN mechanism, replication policy, provisioning
for updates), and every number it reports has a worst-case analysis
behind it. Its distance to the optimum is measured with an exact
formulation on instances of up to 300 task-chains (\cref{sec:rq1}).
